\documentclass[11pt]{article}

\usepackage[margin=1in]{geometry}
\usepackage{amsmath,amssymb,amsthm,mathtools}
\usepackage[T1]{fontenc}
% portability: use lmodern+microtype where available, fall back gracefully
\IfFileExists{lmodern.sty}{\usepackage{lmodern}\usepackage{microtype}}{\usepackage[expansion=false]{microtype}}
\usepackage{booktabs}
\usepackage[hidelinks]{hyperref}

\newtheorem{theorem}{Theorem}
\newtheorem{lemma}[theorem]{Lemma}
\newtheorem{proposition}[theorem]{Proposition}
\newtheorem{corollary}[theorem]{Corollary}
\newtheorem{remark}[theorem]{Remark}
\newtheorem{example}[theorem]{Example}

\newcommand{\R}{\mathbb{R}}
\newcommand{\N}{\mathbb{N}}
\newcommand{\eps}{\varepsilon}

\title{Stopped Channels and Offset Alignment:\\
De-Conditionalizing the Real-Channel Hypotheses\\
for Mixed-Base Tetration\\[2mm]
\large Paper III of the mixed-base tetration series}
\author{Janis Justus}
\date{July 2026}

\begin{document}
\maketitle

\begin{abstract}
The mixed-base phase law of [1] and its diffeomorphism companion [2] are
conditional on two hypotheses about the regular branch: the real-channel
positivity (R2)/(R$'$) and the peeling-growth hypothesis 2.4.  We remove most
of this conditionality.  (i)~The growth hypothesis holds unconditionally for
every admissible tower base $d>e^{1/e}$, with the explicit gap
$T_d(j{+}1)-T_d(j)\ge(1+\ln\ln d)/\ln d>0$.  (ii)~Positivity of the finite
logarithmic chain holds unconditionally down to an explicit \emph{stop level}
$J_0=\min\{j:\ T_d(j)\ge 2C_0\}$ (the \emph{Stoppregel}), by a self-contained
downward induction that produces, rather than assumes, the channel; the
evaluation and limit machinery of [1,\,2] only ever needs the stopped chain.
(iii)~For the limit objects the correct channel statement is
\emph{offset-aligned}: $E_j=T_b(j+\widetilde H(\theta))>0$ is required only
for $j\ge s$, where $s=\lceil -1-\inf\widetilde H\rceil^{+}$ is computable
a priori from the mean-shift hub; the unshifted reading is false for pairs
with $\mu_{b,d}\le-2+\|P\|_\infty$ --- e.g.\ $(b,d)=(10,2)$ and $(e,1.5)$ ---
while the offset version holds with explicit margins.  A numerical
certification across all $45$ measured pairs is reported.  Together with the
base-differentiability of the Kneser family established by the response-kernel
program, the conditional core of the series reduces to a single
mode-certifiable inequality per pair.
\end{abstract}

\section{Introduction}

The phase law of Paper~I \cite{paper} and its diffeomorphism companion,
Paper~II \cite{unified}, are conditional on two hypotheses about the
regular Kneser/Paulsen--Cowgill branch: positivity of the real logarithmic
channel ((R2) of \cite{paper}, (R$'$) of \cite{unified}) and the
peeling-growth hypothesis (Hyp.~2.4 of \cite{paper}). This paper removes
most of that conditionality, in the form in which the estimates of the
series actually consume it.

The contributions are: (i) the growth hypothesis holds \emph{unconditionally}
for every admissible tower base $d>e^{1/e}$, with the explicit gap
$d^{\,x}-x\ge(1+\ln\ln d)/\ln d>0$ (Section~\ref{sec:growth}); (ii)
positivity of the finite logarithmic chain holds unconditionally down to an
explicit stop level $J_0$, by a self-contained downward induction whose
invariant is at the same time the error control of the evaluation ladder
used throughout the series (the \emph{Stoppregel}, Section~\ref{sec:stop});
(iii) for the limit objects, the correct channel statement is
\emph{offset-aligned}: positivity is required only from an explicit shift
$s=\lceil-1-\inf\widetilde H\rceil^{+}$ onward, computable a priori from
the mean-shift hub, and the unshifted reading is provably false for skew
pairs (Section~\ref{sec:offset}); (iv) the alignment margins are certified
numerically across $45$ measured pairs (Section~\ref{sec:cert}). Together
with the base-differentiability theorem of Paper~V \cite{rh}, the
conditional core of the series reduces to a single mode-certifiable
inequality per pair.

\section{Setting}

We keep the conventions of [1,\,2].  For an admissible base $a>e^{1/e}$,
$T_a$ denotes the regular (Kneser/Paulsen--Cowgill) tetration branch with
$T_a(0)=1$, $T_a(x+1)=a^{T_a(x)}$, strictly increasing from $(-2,\infty)$
onto $\R$, and $A_a=T_a^{-1}=\operatorname{slog}_a$.  For an ordered pair
$(b,d)$ set
\[
\alpha=\log_b d,\qquad L_b=\log_b,\qquad
s_j=\inf_{0\le\theta\le1}T_d(j+\theta)=T_d(j),
\]
\[
C_0=\frac{|L_b(\alpha)|+\ln 2/\ln b}{\alpha},
\qquad
W_k^{(n)}(\theta)=L_b^{\circ k}\bigl(T_d(n+\theta)\bigr).
\]
$\Phi_{b,d}$, $\mu_{b,d}=\int_0^1\Phi_{b,d}$, $P_{b,d}=\Phi_{b,d}-\mu_{b,d}$
and $\widetilde H_{b,d}=\mathrm{id}+\Phi_{b,d}$ are as in [1,\,2], and
$A_k=A_k(b,d)$ denotes the Fourier amplitudes of $P_{b,d}$.

Hypothesis (R$'$) of [2] (the quantitative reading of (R2) in [1]) demands
$W_k^{(n)}>0$ for $0\le k\le n-1$.  Hypothesis 2.4 of [1] demands
$s_j\to\infty$.  Both were left conditional in [1,\,2] (see [2, Rem.~7.4]).

\section{The growth hypothesis is free}\label{sec:growth}

\begin{lemma}[Tower divergence with explicit gap]\label{lem:growth}
Let $d>e^{1/e}$.  Then for all $x\in\R$,
\[
d^{\,x}-x\;\ge\;\delta_0(d):=\frac{1+\ln\ln d}{\ln d}\;>\;0 .
\]
Consequently $s_{j+1}=d^{\,s_j}\ge s_j+\delta_0(d)$, hence $s_j\to\infty$;
by [2, Lemma~2.5] the divergence self-improves to
$s_{J_1+k}\ge 2^{k}s_{J_1}$ and $\sum_j 1/s_j<\infty$.
In particular Hypothesis~2.4 of {\rm[1]} holds unconditionally.
\end{lemma}

\begin{proof}
$x\mapsto d^{\,x}-x$ is convex with critical point
$x_*=-\ln\ln d/\ln d$, where the minimum equals
$d^{\,x_*}-x_*=(1+\ln\ln d)/\ln d$.  For $d>e^{1/e}$ one has
$\ln\ln d>-1$, so the minimum is positive.  The remaining statements are
immediate (the inf in $s_j$ is attained at $\theta=0$ since $T_d$ is
increasing).
\end{proof}

\section{The Stoppregel: unconditional finite-chain positivity}\label{sec:stop}

The following lemma re-reads the tail computation of [2, Lemma~6.3]
\emph{unconditionally}: the downward induction produces the positivity of the
chain instead of assuming it, at the price of stopping $J_0$ levels above the
bottom.

\begin{lemma}[Stopped channel]\label{lem:stoppregel}
Let $b,d>e^{1/e}$ and let
\[
J_0=J_0(b,d):=\min\{\,j\ge0:\ s_j\ge 2C_0\,\}\qquad(\text{finite by
Lemma~\ref{lem:growth}}).
\]
Then for every $n\ge J_0+1$ and $\theta\in[0,1]$ the chain values
$W_k^{(n)}(\theta)$, $0\le k\le n-J_0$, are defined and satisfy, for
$1\le k\le n-J_0$,
\begin{equation}\label{eq:invariant}
W_k^{(n)}(\theta)=\alpha\,T_d(n-k+\theta)\bigl(1+\sigma_k\bigr),
\qquad
\sigma_1=0,\qquad
|\sigma_k|\le \frac{C_0}{s_{\,n-k}}\le\frac12\quad(k\ge2).
\end{equation}
In particular $W_k^{(n)}>0$ for $0\le k\le n-J_0$: hypothesis {\rm(R$'$)}
holds unconditionally for the chain stopped at level $J_0$.
\end{lemma}

\begin{proof}
$W_0=T_d(n+\theta)>0$.  The first peel is exact:
$W_1=L_b\bigl(d^{\,T_d(n-1+\theta)}\bigr)=\alpha T_d(n-1+\theta)$, i.e.\
\eqref{eq:invariant} with $\sigma_1=0$.  Assume \eqref{eq:invariant} for some
$k$ with $j:=n-k\ge J_0+1$.  Then $|\sigma_k|\le\frac12$, so $W_k>0$ and
$L_b(W_k)$ is defined; using
$L_b(T_d(j+\theta))=\alpha T_d(j-1+\theta)$ and additivity of $L_b$ on
positive factors,
\[
W_{k+1}
=L_b(\alpha)+\alpha T_d(j-1+\theta)+L_b(1+\sigma_k)
=\alpha T_d(j-1+\theta)\bigl(1+\sigma_{k+1}\bigr),
\]
\[
\sigma_{k+1}=\frac{L_b(\alpha)+L_b(1+\sigma_k)}{\alpha T_d(j-1+\theta)},
\qquad
|\sigma_{k+1}|\le\frac{|L_b(\alpha)|+\ln2/\ln b}{\alpha\,s_{j-1}}
=\frac{C_0}{s_{j-1}}\le\frac12,
\]
the last step because $j-1\ge J_0$.  The induction proceeds while
$n-k\ge J_0+1$, i.e.\ up to $k=n-J_0$.
\end{proof}

\begin{corollary}[Stopped evaluation and stopped tails]\label{cor:stopped}
Define the stopped tails and the stopped phase by
\[
D_n^{\mathrm{stop}}(\theta):=W_{\,n-J_0}^{(n)}(\theta)
=\alpha T_d(J_0+\theta)(1+\sigma_{n-J_0}),
\qquad
\Phi^{\mathrm{stop}}_n(\theta):=(n-J_0)+A_b\bigl(D_n^{\mathrm{stop}}(\theta)\bigr)-(n+\theta).
\]
These are defined for all $n\ge J_0+1$ and $\theta\in[0,1]$ without any
channel hypothesis.  Wherever the full chain of {\rm[1]} exists, the global
Abel equation gives $A_b(D_n^{\mathrm{stop}})=A_b(D_n)+J_0$, so
$\Phi^{\mathrm{stop}}_n=R_n$; the stopped object extends the definition to all
admissible pairs.  All tail estimates of {\rm[2, \S6]} hold verbatim for the
index range $j\ge J_0$.
\end{corollary}

\begin{remark}
Numerically, Corollary~\ref{cor:stopped} is precisely the evaluation ladder
used throughout the measurement campaign (peel down to a level with
$T_d\asymp10^{4}$, hand the remainder to $\operatorname{slog}_b$): the
Stoppregel is not a workaround but \emph{the} definition, and
\eqref{eq:invariant} doubles as its certified error control, cf.\
[2, Rem.~7.5(ii)].  For the standard pairs $J_0\in\{1,2\}$; even at the edge
pair $(e,1.46)$ one has $J_0=23$ (Table~\ref{tab:cert}).
\end{remark}

\section{Offset alignment of the limit channel}\label{sec:offset}

For the limit objects the channel condition concerns
$E_j(\theta)=T_b\bigl(j+\widetilde H(\theta)\bigr)$, $j\ge1$
(cf.\ [2, (11)--(12)]): $E_j>0$ is what makes the $j$-th logarithm real.
Since $T_b>0$ exactly on $(-1,\infty)$ and $T_b$ is increasing, the correct
statement is an \emph{offset} one.

\begin{lemma}[Offset channel]\label{lem:offset}
Assume the phase-law package {\rm(P)} of {\rm[2]} for $(b,d)$, and let
\[
m_\star:=\inf_{\theta}\widetilde H_{b,d}(\theta)\;\ge\;\mu_{b,d}-\|P_{b,d}\|_\infty
\;\ge\;\mu_{b,d}-\sum_{k\ge1}A_k(b,d),
\qquad
s:=\max\bigl(0,\ \lceil -1-m_\star\rceil\bigr).
\]
Then $E_j(\theta)>0$ for all $j\ge s+1$ and all $\theta$, with margin
\[
E_j(\theta)\ \ge\ T_b\bigl(j-s\,+\,(s+m_\star)\bigr)\ >\ T_b(-1)=0,
\qquad j\ge s+1 .
\]
Moreover $D^{(s)}:=T_b(s+\widetilde H)$ is defined (argument $>-2$ requires
only $s+m_\star>-2$, weaker still), and the entire structure of
{\rm[1,\,2]} --- endpoint identities, derivative product, groupoid and
antisymmetry statements --- holds for the $s$-shifted objects by
shift-equivariance of the tetration equation.
\end{lemma}

\begin{proof}
$j+\widetilde H\ge j-s+(s+m_\star)$ and $s+m_\star>-1$ by the definition of
$s$; monotonicity and positivity of $T_b$ on $(-1,\infty)$ give the claim.
Shift-equivariance: replacing $\widetilde H$ by $s+\widetilde H$ replaces
$D$ by $D^{(s)}=T_b^{\circ}\!$-transported one level per unit shift, and every
identity used in {\rm[1,\,2]} is a consequence of $T(x+1)=a^{T(x)}$, which is
invariant under integer shifts of the height coordinate.
\end{proof}

\begin{example}[The unshifted reading fails for skew pairs]\label{ex:fail}
For $(b,d)=(10,2)$: $\mu=-2.26496\ldots$, $\|P\|_\infty\le4.7\cdot10^{-4}$,
so $\inf\widetilde H\le-2.264<-2$: even $D=T_b(\widetilde H)$ is undefined on
part of the period, and $E_1=T_b(1+\widetilde H)$ takes negative values; with
$s=2$ everything is restored, with margin $s+m_\star+1=0.7346$.
For $(e,1.5)$: $\mu=-9.7763\ldots$, $s=9$, margin $0.2234$.
The measured phase law for these pairs (validated to
$10^{-38}$ antisymmetry in the present campaign) is throughout the phase law
\emph{of the offset objects}.
\end{example}

\begin{remark}[Hub anchoring]
$\mu_{b,d}$ --- hence $s$ --- is computable a priori from the hub column
$m(\cdot)$ via $\mu_{b,d}=m(d)-m(b)+\Delta^{(2)}+\Delta^{(3)}+\dots$
(validated to $\sim10^{-15}$ absolute in this campaign), and
$\|P\|_\infty$ from the first few $A_k$ plus the universal decay law.  A
conversion engine can therefore fix the offset \emph{before} evaluating
anything.
\end{remark}

\section{Numerical certification}\label{sec:cert}

Table~\ref{tab:cert} shows an excerpt of the certification (all $45$ measured
pairs pass; script \texttt{foundation\_cert.py}, data at 30--170 digits).
Margins are $s+m_\star+1$ with the rigorous mode bound for $m_\star$.

\begin{table}[h]
\centering\small
\begin{tabular}{lrrrrrr}
\toprule
pair & $\mu_{b,d}$ & $\sum A_k$ & $s$ & margin & $C_0$ & $J_0$ \\
\midrule
$(e,2)$      & $-1.128404$ & $1.5\cdot10^{-4}$ & $1$ & $0.8715$ & $1.529$ & $2$ \\
$(2,e)$      & $+1.128404$ & $1.5\cdot10^{-4}$ & $0$ & $2.1283$ & $1.060$ & $1$ \\
$(10,2)$     & $-2.264961$ & $4.7\cdot10^{-4}$ & $2$ & $0.7346$ & $2.732$ & $3$ \\
$(2,10)$     & $+2.264961$ & $4.6\cdot10^{-4}$ & $0$ & $3.2645$ & $0.822$ & $1$ \\
$(3,5)$      & $+0.578835$ & $1.5\cdot10^{-4}$ & $0$ & $1.5787$ & $0.668$ & $1$ \\
$(e,10)$     & $+1.136557$ & $3.2\cdot10^{-4}$ & $0$ & $2.1362$ & $0.663$ & $1$ \\
$(e,1.7)$    & $-2.879554$ & $2.5\cdot10^{-4}$ & $2$ & $0.1202$ & $2.500$ & $4$ \\
$(e,1.5)$    & $-9.776277$ & $3.6\cdot10^{-4}$ & $9$ & $0.2234$ & $3.936$ & $11$ \\
$(e,1.46)$   & $-21.94334$ & $3.9\cdot10^{-4}$ & $21$ & $0.0563$ & $4.399$ & $23$ \\
$(e,100)$    & $+1.600615$ & $5.7\cdot10^{-4}$ & $0$ & $2.6000$ & $0.482$ & $0$ \\
\bottomrule
\end{tabular}
\caption{Certification excerpt.  The margin is governed by the fractional
part of $\mu$: near-integer $\mu$ produces thin margins (cf.\ $(e,1.46)$),
a structure that also appears to organize the measured singularity abscissae
$x^*(b,d)$ of the strip analysis.}
\label{tab:cert}
\end{table}

\section{What remains, and the differentiability building block}

With Lemmas~\ref{lem:growth}--\ref{lem:offset} the conditional content of
[1,\,2] reduces to the convergence statements of the phase law itself
(Theorem~5.6/5.7 of [1]) formulated for the stopped chain --- estimates that
live entirely in the index range $j\ge J_0$ where \eqref{eq:invariant} is
unconditional --- plus one inequality per pair (the mode-certified
$m_\star$-bound of Lemma~\ref{lem:offset}), which the hub decomposition makes
checkable a priori.

Independently, the response-kernel program of the present campaign
establishes numerically, to $24{+}$ digits and by two independent routes
(complex-step differentiation of the Kneser engine through Paulsen
holomorphy; $\eps$-stencils of the base family), that the Kneser family is
differentiable in the base parameter $\beta=\ln\ln b$, with kernel
$V=\partial_\beta A_b$ characterized by the selection principle
(cohomological equation, normalization $V(1)=0$, prescribed fixed-point
singularity with $dL/d\beta=\ln b\,L^2/(1-\lambda)$, and reality).  This is
the second building block Remark~7.4 of [2] asks for; the constructive
(engine-free) realization of the selection principle is the subject of the
Riemann--Hilbert program documented in the research log.


\section*{Authorship and AI Contribution Statement}
This paper is part of a five-paper series produced in a directed human--AI
collaboration. The author conceived the mixed-base phase program, built and
maintains the high-precision tetration engine used for every computation in
the series (a fork of Levenstein's \texttt{fatou.gp}; engine, data and
scripts at \url{https://github.com/Lightrunnerwastaken/gp-tetration}), and directed the research: posing the questions,
setting priorities between rounds, auditing results, and deciding which
claims met the acceptance gates. Substantial parts of the mathematical
development and of the numerical work were produced by AI systems under
that direction: Anthropic's Claude in a theory and review role, an
autonomous research agent operating the numerical pipeline, and OpenAI's
ChatGPT, which contributed to the original versions of Papers I and II and
to manuscript review. In line with prevailing journal policies, AI systems
are not listed as authors; the author assumes full and sole responsibility
for all statements and proofs, and readers should weigh this provenance in
their assessment. The five papers of the series are permanently archived on Zenodo under the concept DOIs \url{https://doi.org/10.5281/zenodo.21346803} (I), \url{https://doi.org/10.5281/zenodo.21361967} (II), \url{https://doi.org/10.5281/zenodo.21362206} (III), \url{https://doi.org/10.5281/zenodo.21363593} (IV), and \url{https://doi.org/10.5281/zenodo.21363809} (V).

\begin{thebibliography}{9}
\bibitem{paper} J.~Justus, \emph{A Phase Law for Mixed-Base Tetration: Logarithmic Tails, Cocycle Structure, and Singular Asymptotics}, Paper~I of the mixed-base tetration series, preprint, July 2026. \url{https://github.com/Lightrunnerwastaken/gp-tetration}. DOI: \url{https://doi.org/10.5281/zenodo.21346803}.
\bibitem{unified} J.~Justus, \emph{The Mixed-Base Tetration Phase Map Is a Circle Diffeomorphism}, Paper~II of the mixed-base tetration series, preprint, July 2026. \url{https://github.com/Lightrunnerwastaken/gp-tetration}. DOI: \url{https://doi.org/10.5281/zenodo.21361967}.
\bibitem{rh} J.~Justus, \emph{The Base-Derivative of the Kneser Family: Selection Principle, Uniqueness, and Existence via a Riemann--Hilbert Problem}, Paper~V of the mixed-base tetration series, preprint, July 2026. \url{https://github.com/Lightrunnerwastaken/gp-tetration}. DOI: \url{https://doi.org/10.5281/zenodo.21363809}.
\bibitem{kneser} H.~Kneser, \emph{Reelle analytische L\"osungen der Gleichung
$\varphi(\varphi(x))=e^x$}, J.~reine angew.~Math.\ 187 (1950).
\bibitem{pc} W.~Paulsen, S.~Cowgill, \emph{Solving $F(z+1)=b^{F(z)}$ in the
complex plane}, Adv.\ Comput.\ Math.\ 43 (2017).
\bibitem{p19} W.~Paulsen, \emph{Tetration for complex bases}, Adv.\ Comput.\
Math.\ 45 (2019).
\end{thebibliography}

\end{document}
